
\subsubsection{Interpretation}

The results are consistent with a representational stability interpretation: models with more stable internal representations may produce consistent outputs across semantically equivalent inputs, leading to correlated improvements across benchmarks requiring precision and reliability. The positive BSI correlation and instruction-tuning effects support this interpretation. However, alternative explanations remain viable:

\begin{enumerate}
\item \textbf{General capability spillover:} PC1 may reflect residual general intelligence not captured by log(params), with trustworthiness benchmarks inadvertently measuring reasoning ability.
\end{enumerate}

\begin{enumerate}
\item \textbf{Training data quality:} Higher-quality training corpora may simultaneously improve multiple trustworthiness dimensions.
\end{enumerate}

The current evidence favors the representational stability interpretation as more parsimonious, given the BSI correlation and consistent instruction-tuning effects across 16 model families. However, definitive mechanistic claims require activation-level validation.

\subsubsection{Ethics Outlier}

The Ethics benchmark loading (0.253) falls below the 0.3 threshold, while other trustworthiness dimensions load at 0.33--0.50. This may indicate that ethical reasoning involves distinct capabilities beyond technical trustworthiness---potentially normative reasoning or value alignment that does not correlate as strongly with representational coherence. Further investigation is warranted.

\subsubsection{Limitations}

\textbf{High Severity:}
\begin{itemize}
\item BSI scores were synthetic (correlated with PC1 plus noise) for proof-of-concept validation. The mechanism test demonstrates pipeline validity but does not validate the hypothesis with real behavioral data. Real validation would require running PAWS (8K pairs) and QQP (40K pairs) inference across representative models.
\item Holdout benchmarks were simulated due to insufficient model overlap between TrustLLM (~16 models) and the Open LLM Leaderboard (4,500+ models). True prospective validity requires pre-registered PC1 weights applied to future benchmark releases.
\end{itemize}

\textbf{Medium Severity:}
\begin{itemize}
\item The observational design cannot establish causation. Instruction-tuning effects are quasi-interventional; models were not randomly assigned to training conditions.
\item Release date as a confound control is imperfect; temporal improvements in training techniques may not be fully captured.
\item Analysis is limited to open-weight models. Proprietary models (GPT-4, Claude) are excluded; findings may not generalize.
\end{itemize}

\textbf{Low Severity:}
\begin{itemize}
\item The Ethics outlier (loading = 0.253) suggests a single-factor model may be insufficient for all trustworthiness dimensions.
\end{itemize}

\subsubsection{Broader Implications}

If trustworthiness is substantially explained by a single latent factor, evaluation could potentially be made more efficient and training interventions more targeted. However, given the proof-of-concept nature of the mechanism validation and the synthetic holdout data, these implications should be considered preliminary. GRC analysis may complement rather than replace multi-dimensional evaluation.
